\subsection{The output of the application, panel by panel}
The trigger of Chapter 3 is shown in the application by two layers. \emph{Sensor and trigger} follows the suction at the sensor day by day through 20 seasons and marks the days beyond the trigger, and \emph{Trigger law} shows the depletion of the root zone against the trigger and the trade-off between a drier and a wetter trigger.

The line is the daily suction at the sensor depth, the dashed line is the trigger suction (here 471 cm for 50 per cent depletion) and the vertical marks are the irrigations. The bars below count the days beyond the trigger in each year.

\appfig{v1_0}{0.96}{Sensor suction, every day from 2019 to 2039, with the trigger (dashed) and the irrigation events.}

The season is replayed day by day; the text gives the days beyond the trigger, the mean suction and the irrigation and yield of the run, which are the quantities that define the reach of a trigger.

\appfig{v1_1}{0.96}{One season at the sensor: 81 of 183 days beyond 471 cm and a mean suction of 378 cm, in the irrigated run.}

The picture on the left is the root zone with the depletion $f$ coloured; the graph shows the suction and the depletion of the season against the trigger line; and the small drawings give the physical meaning of the numbers, namely that at a trigger of $f=0.50$ the suction is 471 cm in this soil.

\appfig{n1_0}{0.96}{Root-zone depletion and the trigger law: for the soils of the ensemble, the water held between field capacity and wilting, and the trigger at which irrigation fires.}

The table compares two neighbouring triggers on the yield, the irrigation, the drainage and the nitrogen leached, so that the cost of irrigating slightly earlier can be read directly. It is the practical use of the dose law of this chapter.

\appfig{n1_1}{0.96}{The core trade-off: trigger at $f=0.50$ or $f=0.45$?}

